\documentclass[11pt]{article}
\usepackage[margin=0.85in]{geometry}
\usepackage{amsmath,amssymb,booktabs,longtable,array,xcolor}
\usepackage[colorlinks,linkcolor=blue!45!black,urlcolor=blue!45!black]{hyperref}

\newcommand{\M}{\mathcal M}
\definecolor{cut}{RGB}{150,25,20}
\definecolor{put}{RGB}{20,100,55}
\definecolor{hd}{RGB}{45,45,65}

\newcommand{\LOC}[2]{\medskip\noindent\textcolor{hd}{\rule{\textwidth}{0.6pt}}\par
  \noindent\textbf{p.#1, line #2}\par\smallskip}
\newcommand{\NEW}[1]{\medskip\noindent\textcolor{hd}{\rule{\textwidth}{0.6pt}}\par
  \noindent\textbf{#1}\par\smallskip}
\newcommand{\OLD}[1]{\noindent\textcolor{cut}{\textbf{Now:}}\ \begingroup\itshape #1\endgroup\par\smallskip}
\newcommand{\FIX}[1]{\noindent\textcolor{put}{\textbf{Change to:}}\ \begingroup\itshape #1\endgroup\par\smallskip}
\newcommand{\NOTE}[1]{\noindent\textbf{Reason:}\ #1\par}

\setlength{\parindent}{0pt}
\setlength{\parskip}{4pt}

\title{Revision plan}
\author{}
\date{28 August 2026}

\begin{document}
\maketitle

This covers everything: the story, the title, the new results, the citations, the figures, and
the line edits. Page and line numbers are the margin numbers of the 28 August draft.

% ==================================================================
\section{The story to tell}

The paper currently argues that transport can be harmful and that conditioning on the source
makes things worse. Both are true, but as a framing this has three problems. It reads as a
warning rather than a finding. It contradicts two methods that work, one of which was written by
an invited speaker at the venue. And it measures a single axis, cosine similarity, which the
paper itself spends a section explaining you should not trust.

There is a better story available, and it needs no new experiments.

Two things can happen when you transport an informative source. You can trade: give up pointwise
closeness to the particular target and gain distributional realism. This is the
distortion-perception tradeoff, and it has been known since 2018. Or you can lose on both axes at
once, which is not a trade and has no reason to happen unless something is structurally wrong.

The paper's own numbers separate these two cases cleanly, and nobody has separated them before.
Increasing integration accuracy trades: cosine falls from $0.3541$ to $0.3172$ while polar $W_1$
falls from $0.4623$ to $0.0321$. Stable Diffusion at $8\times$ degradation trades: cosine gain
$-0.055$, energy gain $+0.054$. Truncating the early interval trades, badly: cosine rises to
$0.2997$ while $W_1$ rises to $0.6939$, which is ninety per cent of the way back to doing nothing.

Source conditioning does not trade. $\M_2$ is worse than $\M_1$ on cosine, worse on pairwise
cosine, and worse on $W_1$, in all twelve learning-rate by width cells, at both resolutions of
the super-resolution experiment, and under both the velocity and the endpoint parameterisation.
That is the finding.

And it explains two methods rather than contradicting them. PMRF starts from the posterior mean,
the most informative source available, and transports it, and it wins. It wins because its
velocity field is not conditioned on the source: the source only sets the initial condition. In
this paper's terms PMRF is $\M_1$. Augmented Bridge Matching supplies the initial sample to the
velocity field, explicitly trading the Markov property to recover the training coupling, and it
wins on that objective. In this paper's terms it is $\M_2$. The $\M_1$ against $\M_2$ comparison
is the controlled measurement of exactly the choice that separates them.

So the paper is not saying do not transport. It is saying that where the source enters matters,
that the cost of getting it wrong is measurable, and that in one setting the boundary can be
computed in closed form.

% ==================================================================
\section{Title and abstract}

\NEW{Title, p.1}
\OLD{When Conditional Flow Matching Should Do Nothing: Informative Endpoint Coupling and
Source-Conditioning Failure}
\FIX{Transport Moves You Along the Curve; Conditioning Moves You Off It}
\NOTE{The proposed title is good and I would use it. One caveat so you use it knowingly. The
second half is exactly right: $\M_2$ is worse on every axis measured, which is off the curve. The
first half is right for how \emph{much} transport you apply (integration steps, SDEdit strength),
where cosine and distribution genuinely move in opposite directions. It is not right for
\emph{whether} to transport at all: above $\beta^\star$ the flow beats doing nothing on cosine
\emph{and} on $W_1$ ($0.048$ against $0.767$), which is strictly better, not a trade. If you want
the title to be literally defensible everywhere, add a subtitle that scopes it, for example
``Transport Moves You Along the Curve; Conditioning Moves You Off It: Informative Sources in
Conditional Flow Matching''. If you are comfortable with a title that is a slogan rather than a
theorem, use it as is. Most people will.}

\NEW{Abstract, p.1, lines 1--21}
\OLD{The current abstract leads with ``Conditional flow matching typically treats the source as
generic initialization'' and ends with ``informative sources change both when conditional flow
matching should transport and how source information should be supplied to the dynamics''.}
\FIX{Draft replacement, to be tightened for length:

``When the source of a conditional flow is already informative about the target, two questions
separate that coincide under noise initialization: whether to transport it, and whether to tell
the velocity field what it was. We study both in a spherical model where the endpoint coupling
can be varied at fixed conditional marginals, and we give a closed-form landing law and an
explicit threshold locating the first question. For the second we show that supplying the source
to the field makes the flow-matching target exactly recoverable along training trajectories, at
the cost of a recovery map whose sensitivity diverges as $1/t$; the same divergence holds on the
geodesic interpolant actually used. Empirically the two questions behave differently. Applying
more transport trades pointwise alignment for distributional fidelity, the tradeoff of Blau and
Michaeli, and we locate it exactly. Conditioning the field on the source does not trade: it is
worse on target alignment, on conditional diversity, and on distributional distance
simultaneously, across twelve learning-rate and width settings, at three super-resolution
resolutions, and under both velocity and endpoint parameterization. This distinction explains why
posterior-mean rectified flow, which transports a maximally informative source without
conditioning on it, succeeds, while augmentation that conditions on the source succeeds only for
objectives that reward coupling recovery. Experiments with OpenCLIP embeddings, super-resolution,
a masked-diffusion analogue, and a pretrained Stable Diffusion image-to-image model test the
picture beyond the controlled setting.''}
\NOTE{The current abstract contains no word that connects to the venue's framing and no
acknowledgement of the literature it sits in. This version names the tradeoff, names the two
methods it reconciles, and reports both axes.}

% ==================================================================
\section{New results to insert}

Six completed experiments are absent from the draft. Two of them change what the paper can claim.

\NEW{New subsection in Appendix B: the parameterisation control}
\FIX{Insert the text of \texttt{param\_section.tex}, first block.}
\NOTE{This is the answer to the first question anyone will ask, and it was open until this week.
The gap survives endpoint prediction: at $\eta=10^{-4}$ it retains $96\%$ of its size at $80$
standard errors, at $\eta=10^{-3}$ it retains $70\%$, and averaged over rates $78\%$. The
diversity collapse survives, and $\M_2$ still attains the lower training loss under endpoint
prediction. Report the qualification too: roughly $30\%$ of the extra gap induced by the larger
step size is parameterisation-specific, even though the gap itself is not.}

\NEW{Section 8 Limitations, p.8, lines 303--305, and a new appendix subsection}
\FIX{A discrete masked-diffusion analogue was run over six source-informativeness levels with
four arms. Conditioning on the source while starting from it \emph{helps} there, by $+0.0107$ on
average and at five of six levels, where in the continuous model it costs about $0.10$.}
\NOTE{The limitation paragraph currently says whether the mechanism persists under stochastic
dynamics ``remains open''. It is not open. It was tested and the sign reversed. Reporting a
boundary you found yourself is much better than having a reviewer find it. Full table in
\texttt{review\_aug28}, section 4.3.}

\NEW{Appendix A.2.1, after p.11 line 391}
\FIX{Add the criterion verified across dimensions $2$ to $64$: systems satisfying both hypotheses
have worst ninetieth-percentile drift $4.06\times10^{-8}$, systems violating a hypothesis at
relative size $0.1$ have smallest median drift $9.78\times10^{-3}$, a separation of
$2.41\times10^{5}$.}
\NOTE{The paper currently supports Proposition 1 with a $180$-point power-law fit. This is a much
stronger statement and it costs two sentences.}

\NEW{After Proposition 1, p.3, lines 119--123}
\FIX{Add the four-arm test isolating the two hypotheses. Symmetric and in span, commuting
marginals: $6.34\times10^{-5}$. Symmetric and in span, non-commuting: $6.34\times10^{-5}$.
Symmetric and off span: $7.62\times10^{-3}$. Skew and in span: $9.52\times10^{-1}$.}
\NOTE{Two things the paper asserts without evidence follow immediately. The remark at p.11 lines
381--383, that invariance does not need commuting marginals, is confirmed to five digits by the
first two rows. And the $6.585$ asymmetry appears here directly as a factor of $125$.}

\NEW{Appendix B.2, near p.17 line 585}
\FIX{Add the standard-parameterisation control: at width $2048$ the gap is $0.0998$ under
standard parameterisation and $0.0988$ under maximal-update.}
\NOTE{The capacity conclusion currently rests on one parameterisation. This is the obvious
question and the answer is already computed.}

\NEW{Appendix B.3, near p.20 lines 658--664}
\FIX{Add the integrator anchor test. At $\alpha=0$ and $\sigma_{\mathrm{aux}}=0$ both auxiliary
families use the same trained model, terminal losses $4.9116\times10^{-4}$ and
$4.9046\times10^{-4}$, yet reported distributional errors of $0.0782$ and $0.0367$ when sampled
with different integrators. Under one integrator the anchor agrees and the families trace one
curve, cross-to-within ratio $1.40$ against $3.44$.}
\NOTE{Without this, a reader comparing the two families is comparing two integrators.}

\NEW{Appendix B.8 and the transfer paragraph, p.28 lines 930--941}
\FIX{Add the six-factor super-resolution sweep and the Stable Diffusion distributional rescore.
The rank bound is hit exactly at all six factors, taking the total to nine exact agreements. The
trained flow beats doing nothing at every factor with the gain \emph{falling} as the source
degrades: $3.81$, $2.26$, $1.97$, $1.29$, $1.02$, $0.59$ PSNR. The rescore of the pretrained
model, four samples per image:

\begin{center}
\begin{tabular}{@{}rrrrr@{}}
\toprule
degradation & source align & best cosine gain & best energy gain & pairwise cosine at $s=0.8$\\
\midrule
$2\times$  & $0.997$ & $-0.189$ & $-0.174$ & $0.742$\\
$4\times$  & $0.853$ & $-0.061$ & $+0.008$ & $0.736$\\
$8\times$  & $0.733$ & $-0.055$ & $+0.054$ & $0.719$\\
$16\times$ & $0.430$ & $+0.120$ & $+0.392$ & $0.706$\\
\bottomrule
\end{tabular}
\end{center}
Energy null floor $0.011$; the retained source has pairwise cosine $1.000$ by construction.}
\NOTE{Three things follow. At $2\times$ the do-nothing regime holds on both instruments, so the
headline no longer rests on the metric you disclaim. At $4\times$ the energy gain of $+0.008$ is
below the null floor and must not be read as a reversal. At $8\times$ the two instruments
genuinely invert, which is a fourth instance of the paper's own tradeoff, in a production model.
The contrast with the factor sweep is the useful part: a flow trained for its exact degradation
beats doing nothing everywhere, a pretrained one applied off the shelf does not.}

\NEW{Appendix B.7, near p.27 lines 885--893}
\FIX{Add the held-out label probe: $0.058$ accuracy against a chance level of $0.10$ on a
disjoint split, so the component of the source orthogonal to the condition carries no label
information.}
\NOTE{Phrase it as ``does not extract label information'', not as an accuracy figure. The value
sits about five standard errors below chance, which means the probe anti-generalises rather than
that it is informative.}

% ==================================================================
\section{Mathematics to insert}

\NEW{Theorem 2 and its proof, p.4 lines 153--157 and p.14 lines 494--502}
\OLD{Therefore, $A<\cos\jmath$ is a sufficient condition for the existence of an interior
flow-versus-do-nothing threshold. \emph{Plus the TODO asking for the converse.}}
\FIX{A transition exists for every $A\in(0,1)$ and $\jmath\in(0,\pi/2)$. Substituting
$\beta=\sin\varphi$ and $\psi=\arccos(A\cos\varphi)$ gives
$\Delta/A=\cos(\psi-\varphi)-\cos(\varphi-\jmath)$. At $\varphi=\jmath$ the second term is
$\cos 0=1$, its maximum, while $\psi(\jmath)>\jmath$ strictly for $A<1$, so
$\Delta(\sin\jmath)<0$ always and a root lies in $(\sin\jmath,1)$ unconditionally. What
$A<\cos\jmath$ characterises is uniqueness: $\psi'\le A$ makes $h=2\varphi-\jmath-\psi$ strictly
increasing, so there is exactly one root above $\sin\jmath$, and below it the condition reduces
to $A\cos\varphi>\cos\jmath$, monotone, so a second root appears exactly when $A>\cos\jmath$.}
\NOTE{Verified on $120$ cells with no exceptions. Full proof in \texttt{threshold.tex}. This also
fills three cells of Table B1 and takes the comparison from $27$ to $30$.}

\NEW{Appendix A.6, p.15 lines 521--545}
\FIX{Add the geodesic version. On the interpolant actually used, the recovery determinant
collapses to $-\sin((1-\alpha)t\theta)/\sin\theta$ and the sensitivity is
$\sin((1-\alpha t)\theta)/\sin((1-\alpha)t\theta)$, which diverges as $1/t$ exactly as in the
flat case, differing by the bounded factor $\sin\theta/\theta$.}
\NOTE{Propositions 2 and 3 are derived for the linear interpolant while every experiment uses the
geodesic, and Section 8 flags this as open. It is not. Verified against finite differences in
dimension $64$ to $2.7\times10^{-10}$.}

\NEW{Appendix A.2.1, p.12 lines 426--435}
\FIX{Delete the three TODO blocks. Insert the first-order derivation from \texttt{ratio.tex}, and
correct the explanation of the slope: the ratio slope is identically the skew slope minus the
span slope. Over the reported window the span exponent is $0.573$ and the skew exponent $-0.159$,
giving $-0.732$; a decade lower they are $0.933$ and $-0.019$, giving $-0.952$.}
\NOTE{The shortfall from $-1$ is the span response being sublinear over a window that is not yet
asymptotic. It is not the decline of the skew constant, which pushes the slope toward $-1$, not
away from it. If the paper says otherwise the arithmetic contradicts it.}

% ==================================================================
\section{Citations}

The draft cites nine works. Your own prior-art audit covers sixteen, and five of those are
missing, including one the audit marks ``Same'' with the action written as ``Cite''. Four more
came out of the literature check. Adding these fixes the thin reference list in one pass.

\begin{longtable}{@{}p{0.30\textwidth}p{0.64\textwidth}@{}}
\toprule
Work & Where it goes and what to say\\
\midrule
\endhead
Blau and Michaeli, CVPR 2018 &
Section 2 and the reference-ladder discussion at p.6 lines 207--214. Table C1's ordering, with a
perfect sampler at $0.137$ scoring below doing nothing at $0.264$, is their theorem reproduced on
a sphere. Claim the closed-form quantification, not the phenomenon. Your audit already lists this
as ``Same''.\\
\addlinespace
Ohayon, Michaeli, Elad, PMRF, arXiv:2410.00418 &
Section 2 and Appendix B.3. PMRF is $\M_1$ with a noised initialization. It also adds Gaussian
noise to the posterior mean and gives the reason: source and target lie on manifolds of different
dimension and the noise alleviates the singularity of a deterministic map between them. That is
your Corollary 4.21 and your rank table. Cite it as convergent evidence, and state that your
$\sigma_{\mathrm{aux}}$ perturbs the auxiliary input while the sampler still starts at the
unperturbed source, which is a different intervention.\\
\addlinespace
Freirich et al., distortion-perception function &
Section 2, alongside Blau and Michaeli, for the closed-form treatment.\\
\addlinespace
Kim et al., Better Source Better Flow, arXiv:2602.05951 &
Section 2. They learn a condition-dependent source distribution and identify distributional
collapse plus a variance-regularization fix. Distinguish in one sentence: they regularize the
source distribution, you noise the auxiliary input with the source held fixed.\\
\addlinespace
Albergo et al., stochastic interpolants &
Appendix A.6, as the reference PMRF uses for the low-dimensional-manifold singularity.\\
\addlinespace
Delbracio and Milanfar, InDI &
Section 2, for iterative restoration from a degraded observation. Already in your audit.\\
\addlinespace
Tsimpos et al. &
Section 2 and Appendix A.2. Already in your audit, and their open problem about intermediate
regimes between independence and determinism is the one Theorem 4.11 answers a case of.\\
\addlinespace
Zhang et al., Bridge Graphical Models &
Appendix on the loss floor. They study the same Markovization gap and report the coupling moving
it by two orders of magnitude. Your addition is that the gap is decoupled from, and ordered
against, the endpoint map. State this explicitly or a reader will take their result as a
contradiction.\\
\addlinespace
Pooladian et al., Multisample Flow Matching &
Section 2, coupling design.\\
\bottomrule
\end{longtable}

\textbf{One warning.} Nine sources in your audit are marked as unverified secondary summaries,
with a standing note that each must be confirmed before submission. The load-bearing one is Wang
et al., the prior for the $1/t$ mechanism in Section 6 and Appendix A.6. If it proves more than
the summary suggests, that section narrows. This is an hour of reading, not an experiment.

% ==================================================================
\section{Figures}

The paper has no images despite two image experiments. Both panels now exist.

\begin{longtable}{@{}p{0.22\textwidth}p{0.72\textwidth}@{}}
\toprule
Figure & Action\\
\midrule
\endhead
Figure 1, p.6 line 206 &
Unreadable at three-across. Use the three separate panels, or give the composite full page width.
Panel (a) must label the green points ``population field (measured)'', not ``flow'': they come
from the fitted oracle, which is why they sit near $0.335$ rather than $\M_1$'s $0.3184$.\\
\addlinespace
Figure 2, p.17 line 574 &
Plots the $j=0$ points that Table B1's own caption calls formal only, and the three cells with
$A\ge\cos\jmath$, which is why isolated points appear at $A=0.9$ with no curve. Adopt the
corrected Theorem 2 and plot all $30$ with their analytic upper roots, so every point has a
curve.\\
\addlinespace
Figure 3, p.18 line 591 &
Legend entries are malformed, reading ``source-conditioned $\eta$:0.000 1''. Regenerate.\\
\addlinespace
Figures 1 and 5 &
The same figure appears twice, p.6 and p.24. Fixed by deleting Appendix C.\\
\addlinespace
New: super-resolution panel &
Rows are degraded input, regression, $\M_1$, $\M_2$, target. $\M_2$ is visibly smoother than
$\M_1$ while scoring higher, and the regression row is sharpest of all while scoring worst. This
makes the central claim visible instead of asserted. Put it in the main text if space allows.\\
\addlinespace
New: image-to-image panel &
Rows are degradation levels, columns are clean target, retained source, and outputs at three
strengths. At $2\times$ and strength $0.8$ the model has replaced the input with a different
image, which is why alignment collapses. Caption should note that at $16\times$ the output is a
plausible image but not a recovery of the original, so both metrics may be rewarding realism.\\
\addlinespace
New: two-metric panel &
Cosine gain and energy gain side by side across strengths, with the null floor shaded. This is
the figure that carries the two-axis story.\\
\bottomrule
\end{longtable}

% ==================================================================
\section{Line edits}

These are corrections to text that is already there. The full list with exact replacement wording
is in \texttt{errata.tex}; the ones that change a claim are repeated here.

\begin{longtable}{@{}p{0.14\textwidth}p{0.38\textwidth}p{0.42\textwidth}@{}}
\toprule
Location & Now & Change to\\
\midrule
\endhead
p.7, 242--243 & cosine from $0.350$ to $0.319$ & $0.3541$ to $0.3172$\\
p.7, 252--253 & alignment $0.50$ to $0.69$, prototype $0.79$ & $0.5517$ to $0.7586$, prototype
$0.7930$, random-pair floor $0.4815$\\
p.7, 264 and p.28, 922 & At both completed resolutions & At all three resolutions. At
$96\times96$ the two models score $23.1984$ and $23.1984$, identical, while energy is $1.888$
against $2.222$\\
p.7, 227--228 & gap $0.0985$ at width 256 and $0.1040$ at 2048 & positive in all twelve cells,
$0.0419$ to $0.1158$, minimum $16$ standard errors (Table B2)\\
p.22, 719--721 & $+0.157$ and $-0.505$ & $+0.0770_{\pm0.0004}$ and $-0.4905_{\pm0.0039}$\\
p.22, 722 & a stray \texttt{:contentReference[\ldots]} string & delete\\
p.22, 749--751 & factorization discrepancy is $0.62$ & delete. No such quantity is computed; the
stage was re-run and no column equals $0.62$. Nearest is $\hat\beta=0.6010$, already named\\
p.24, 790--791 & $0.350$ and $0.319$ & $0.3541$ and $0.3172$\\
p.27, 881--882 & $0.2195\to0.3122$, $\M_1=0.3195$ & $0.2149\to0.3105$, $\M_1=0.3180$\\
p.16, 571--573 & maximum absolute error $3.3\times10^{-2}$ & add: the error is one-sided,
measured exceeds analytic in $22$ of $27$ cells, mean signed error $+0.0099$, sign test
$p=1.5\times10^{-3}$\\
p.6, 207--209 & $0.322$ (learned flow) & $0.3184$, matching line 217 and Table C4\\
p.6, 209 & $0.137$ (perfect sampler) & say which reference: $A^2=0.1369$ for a draw sharing the
posterior mean, against $0.0500$ for a draw given the condition alone\\
p.18, Table B3 & the earlier width sweep & delete; fold its loss and distributional trend into
Table B2's discussion\\
\bottomrule
\end{longtable}

\textbf{Add to the evaluation protocol.} The distributional report fixes its evaluation seed and
the fidelity report does not, so scoring one identical trained network twice gives cosines
differing by about $0.002$. State this once. It sets a floor below which differences are not
meaningful, and it resolves several places where the same model appears with slightly different
numbers without editing any of them. All gaps claimed in the paper are between $0.04$ and $0.12$.

% ==================================================================
\section{Structure}

\NEW{Delete Appendix C, p.24 lines 769--959}
\NOTE{Appendix C repeats Appendix B in full. The threshold, the reference ladder, guidance,
capacity, CLIP and super-resolution each appear twice with different numbers, and the note at
p.24 lines 770--773 already says C should go. Move Tables C1 to C9 and Figure 5 into Appendix B,
keeping the C version where it is newer: C6 supersedes the guidance text in B.6, C7 and C8 fill
B.8, C1 belongs in B.4. This removes most of the internal contradictions on its own, and it is
the largest single reduction in similarity-check exposure, since the tool compares the document
against itself.}

\NEW{Nine TODO blocks}
\NOTE{p.12 lines 427--435, p.14 lines 498--502, p.19 lines 610--616, p.22 lines 711--714, p.22
lines 735--739, p.23 lines 754--757, p.23 lines 765--768, p.28 lines 927--929, p.29 lines
942--944. All nine are answered in \texttt{todo12.tex}. None needs new computation.}

\NEW{NeurIPS checklist, p.30 to p.36}
\NOTE{All sixteen answers are \texttt{[TODO]} and the instruction block that the template says to
delete is still at p.30 lines 960--988. The checklist states that papers without it are desk
rejected. Answers with justifications are in \texttt{checklist.tex}. Read item 16 on LLM usage
before pasting it; that one is a judgment call and it is yours to make.}

% ==================================================================
\section{Order of work}

\begin{enumerate}
\item Delete Appendix C and merge its tables into B. Most of the contradictions go with it.
\item Add the ten citations and write the two positioning paragraphs, PMRF and ABM. This is the
highest-value day in the whole plan: it is what makes the paper look like it knows its field.
\item Retitle, rewrite the abstract, and rewrite Section 5 and Section 7.1 around the two axes.
\item Apply the line edits and fill the nine TODOs.
\item Insert the new results: parameterisation control, masked diffusion, criterion at scale,
four-arm test, six-factor sweep, image-to-image rescore, resolution $96$.
\item Insert the three new figures and regenerate Figures 1, 2 and 3.
\item Complete the checklist.
\item Verify the nine unread secondary sources, Wang et al. first.
\end{enumerate}

Nothing above needs a GPU. Every experiment is finished.

\end{document}
